\documentclass{article}
\usepackage{amsmath,amssymb}
\usepackage[margin=1in]{geometry}
\begin{document}
% BEGIN ACCEPTANCE UNCERTAINTY EXPLANATION
Acceptance probabilities fluctuate even under a fixed kernel. A temporal block
below the practical band and another above it therefore do not by themselves
show a change in the chain's acceptance distribution. The default raw crossing
screen is a heuristic, and its separate four-chain compatibility interval does
not calibrate the crossing decision. Independent random seeds also do not make
four differently initialized chain means identically distributed.

For a stationary acceptance sequence $a_{it}$ with an applicable central limit
theorem, uncertainty in its mean depends on the long-run variance
$\sigma_i^2=\gamma_i(0)+2\sum_{k\geq1}\gamma_i(k)$, where $\gamma_i(k)$ is
its lag-$k$ autocovariance. Independence between $m$ chains of length $n$
then gives the approximation
\[
 \widehat{\operatorname{Var}}(\bar a)
 = \frac{1}{m^2 n}\sum_{i=1}^{m}\widehat\sigma_i^2.
\]
This follows by adding the variances of the independent chain means and
dividing by $m^2$. Batch means and lugsail estimate these long-run variances;
the formulas and assumptions are discussed in the diagnostics chapter. A
positive estimate cannot establish adequate batch length, stationarity or
sufficient burn-in.

An experimental acceptance diagnostic now preserves the original tuning
decision while computing these uncertainty estimates separately. It compares
all temporal windows within each chain, so opposing drifts cannot cancel by
averaging across chains. It subtracts simultaneous marginal interval endpoints:
if both means are covered, their difference is covered regardless of their
covariance. A difference interval wholly outside a declared tolerance supports
a material discrepancy; an interval wholly inside supports compatibility for
that contrast; overlap with the tolerance boundary remains unresolved. The
approximate marginal intervals still require calibration, including across
the declared evidence looks and candidate search.

The October 1 validation did not support a new default. Under strongly
persistent stationary simulations, the tested batch rules underestimated
uncertainty and reported false conflicts. Increasing the batch length reduced
some conflicts without establishing adequate coverage. The new diagnostic
therefore has no tuning-artifact authority. Its finite-start model checks and
BGS replays are separate from posterior convergence, and R-hat remains
reporting-only during tuning.

\end{document}
